\begin{theorem}[Theorem~1.3]
Let $\mathbb G$ be a sub-elliptic setting (\noderef{CarnotSetting}) with group law $*$,
homogeneous norm $\|\cdot\|$, step $r$, frame $\sigma=[\sigma^1,\ldots,\sigma^m]$ and
first-layer projection $\pi_1:\mathbb R^{n}\to\mathbb R^{m}$. Assume the cited results of the
literature \emph{at $\mathbb G$} (\noderef{CitedResults}), and also the cited results of the
literature \emph{at the Euclidean setting $\mathbb G_E$ on $\mathbb R^{m}$}
(\noderef{EuclidSetting} at $N:=m$, and \noderef{CitedResults} at that setting). Let
$A:\mathbb R^{m}\to\mathbb R^{m\times m}$ and $H:\mathbb R^{m}\to\mathbb R$ be continuous on
$\mathbb R^{m}\mns\{0\}$, with $A(\xi)$ symmetric and positive semidefinite for $\xi\ne0$,
and satisfy the conditions of strict ellipticity in the gradient direction (E) (1.9) and of
homogeneity along rays (S) (1.10) --- that is, let $(A,H)$ satisfy the structure conditions
of \noderef{StructureConditionsES}. No normalization of $(A,H)$ is assumed. Let
$\Om\subset\mathbb R^{n}$ be a bounded open set with $\mathbb R^{n}\mns\Om\ne\emptyset$, let
$u,v:\mathbb R^{n}\to\mathbb R$ be continuous on $\bar\Om$, and assume
\begin{enumerate}
\item $u$ is a viscosity sub-solution (\noderef{IsViscSubsol}) and $v$ a viscosity
  super-solution (\noderef{IsViscSupersol}) of $\mathcal L\cdot=0$ in $\Om$;
\item $u\le v$ on $\del\Om$.
\end{enumerate}
Let $R\ge0$ satisfy $|u|\le R/2$ and $|v|\le R/2$ on $\bar\Om$. Write
\[
 E:=\argmax_\Om(u-v)=\big\{y\in\Om:\ u-v\ \text{attains its maximum over}\ \bar\Om\ \text{at}\ y\big\}.
\]
Assume finally that at least one of the following three cases holds.
\begin{enumerate}
\item[(i)] There is an open set $N\subseteq\mathbb R^{n}$ with $E\subseteq N$ such that
  either $v$ is $C^{1}$ on $N$ and $\mathrm{hgrad}\ \sigma\ z\ D\ne0$ (\noderef{hgrad})
  whenever $z\in N$ and $D$ is the derivative of $v$ at $z$, or $u$ is $C^{1}$ on $N$ and
  $\mathrm{hgrad}\ \sigma\ z\ D\ne0$ whenever $z\in N$ and $D$ is the derivative of $u$
  at $z$.
\item[(ii)] $u$ is constant on $\bar\Om$, or $v$ is constant on $\bar\Om$.
\item[(iii)] $r=1$; or there are $F,Gf:\mathbb R^{m}\to\mathbb R$, continuous on
  $\pi_1(\bar\Om)$, with $u=F\circ\pi_1$ and $v=Gf\circ\pi_1$.
\end{enumerate}
Then $u\le v$ in $\Om$.
\end{theorem}
\begin{proof}
This is Theorem~1.3. The paper states it as "Let $\Om$, $\mathcal L$, $u$, $v$ be as in
Theorem Theorem~1.2, without the hypothesis on $E_{\eps_k}$", a citation of another theorem's
hypotheses; FILESPEC and the substantiveness rubric require those hypotheses to be written
out, and the statement above does so. Two of them deserve the same justification they were
given on Theorem~1.2.

\emph{The parameter $R$.} The paper defines $R:=2\max\{\|u\|_{L^\infty},\|v\|_{L^\infty}\}$
just before (3.15), whereas \noderef{ReducedDomain} carries $R$ as a parameter. We
therefore quantify over every admissible $R$, i.e. every $R\ge0$ with $|u|,|v|\le R/2$ on
$\bar\Om$. This contains the paper's statement: the paper's $R$ is admissible, since
$\|u\|_{L^\infty}\le\max\{\|u\|_{L^\infty},\|v\|_{L^\infty}\}=R/2$ and likewise for $v$, and
it is the smallest admissible one. It is also not stronger than the paper's statement in any
way the paper does not justify: $R$ enters the argument only through the definition of
$\Om[\eps]$, where a larger $R$ makes $\Om[\eps]$ smaller, and the role of the bound
$|u|\le R/2$ is exactly the one the paper gives it in the attainment step of Lemma~3.8(4),
namely to make $\|z_x\|^{q}\le4\eps\|u\|_{L^\infty}\le2R\eps<(1+4R)\eps$.

\emph{The hypothesis $\mathbb R^{n}\mns\Om\ne\emptyset$.} This is the paper's own
convention, made explicit: (3.15) reads
$\Om[\eps]=\{x\in\Om:\ \inf_{y\in\mathbb R^{n}\mns\Om}\|y*x^{-1}\|^{q}>(1+4R)\eps\}$, an
infimum over $\mathbb R^{n}\mns\Om$, and the paper's $\Om\subset\mathbb R^{n}$ is a proper
subset. It is automatic as soon as $n\ge1$, since no bounded subset of $\mathbb R^{n}$ is
then all of $\mathbb R^{n}$; it excludes only the degenerate setting $n=0$, in which
$\del\Om=\emptyset$ and hypothesis (2) carries no information at all.

\emph{"Constant" in case (ii).} The paper writes "$u$ or $v$ is constant" without naming the
set. We read it on $\bar\Om$. That costs nothing and loses nothing: if $v$ is constant on
$\Om$ then, $v$ being continuous on $\bar\Om$ and $\Om$ dense in $\bar\Om$, $v$ is constant
on $\bar\Om$ as well (and if $\Om=\emptyset$ then $\bar\Om=\emptyset$ and the condition holds
vacuously), so the paper's reading implies ours and the theorem below is at least as strong
as the printed one.

The structure conditions are kept in the form (E) and (S), with no normalization: the printed
theorem is for general $A,H\in C(\mathbb R^{m}\mns\{0\})$ satisfying (1.9) and (1.10), and
(2.3) is an internal step of the proofs, discharged inside the nodes cited below. And
the cited results of the literature are carried as an explicit hypothesis, as pm-0001 and
item 2 of the paper's "Note for the formalization" require, so that the formalized theorem
reads "cited results $\Rightarrow$ Theorem Theorem~1.3".

\emph{Why there are two cited-results hypotheses, and not one.} Three of the five cited
results --- the Bardi--Goffi principles, Pansu's theorem with [DMV13, Prop. 4.8], and parts
(1)--(3) of Lemma~3.8 --- are theorems about Carnot groups, and \noderef{CitedResults} carries
the setting at which they are asserted as a parameter. Asserting them instead at every
instance of \noderef{CarnotSetting} would over-reach: that structure has no dilations and no
stratification, which is precisely what their proofs use, so for an arbitrary instance they
are not known and may fail; a hypothesis package that may be unsatisfiable would make this
theorem vacuous. Item 3 of the paper's note for the formalization licenses building the
structural facts of Section~2.4 into the setting only "provided every Carnot group in the
coordinates of Section Section~2.4 satisfies that definition", which guarantees that every
Carnot group is an instance and says nothing in the converse direction. Case (iii) applies
the cited results a second time, inside the Euclidean setting on $\mathbb R^{m}$ to which it
descends, and a package at $\mathbb G$ does not supply one at that setting; hence the second
hypothesis. It is satisfiable for the same reason the first is: $\mathbb G_E$ is the Carnot
group $(\mathbb R^{m},+)$ with the Euclidean norm, where Pansu's theorem is Rademacher's,
Bardi--Goffi is the classical strong maximum principle, and [CYW07, Prop. 3.3] is the
classical sup-convolution lemma. Neither package is proved here; both are assumed, as item 2
of the note directs.

\medskip
The proof is the three-way case split the paper itself makes ("Only case (i) is obtained
through the hypothesis of Theorem Theorem~1.2; case (ii) follows from the strong maximum
principle alone, and case (iii) from the Euclidean translation argument"). Each case is a
separate node, and the three arguments are genuinely different; the split is the one
contemplated by the note on cases in the substantiveness rubric.

\medskip\noindent\textbf{Case (i).} Let $N$ be as in (i): open, with $E\subseteq N$. The
hypotheses of \noderef{UncondCaseC1} are exactly the ambient hypotheses of the present
statement together with this $N$: the cited results at $\mathbb G$; the structure conditions
on $(A,H)$;
$\Om$ bounded, open, with non-empty complement; $u,v$ continuous on $\bar\Om$; the two
viscosity properties; $u\le v$ on $\del\Om$; $R\ge0$ with $|u|,|v|\le R/2$ on $\bar\Om$; $N$
open; and $E\subseteq N$. That node has two conclusions, one for each orientation of the
$C^{1}$ hypothesis:
\begin{itemize}
\item if $v$ is $C^{1}$ on $N$ and $\mathrm{hgrad}\ \sigma\ z\ D\ne0$ for every $z\in N$ and
  every derivative $D$ of $v$ at $z$, then $u\le v$ in $\Om$;
\item if $u$ is $C^{1}$ on $N$ and $\mathrm{hgrad}\ \sigma\ z\ D\ne0$ for every $z\in N$ and
  every derivative $D$ of $u$ at $z$, then $u\le v$ in $\Om$.
\end{itemize}
Case (i) supplies one of the two antecedents, so one of the two conclusions applies and gives
$u\le v$ in $\Om$.

\medskip\noindent\textbf{Case (ii).} The hypotheses of \noderef{UncondCaseConstant} are the
ambient hypotheses of the present statement, minus $R$ and the bound on $|u|,|v|$, which that
node does not use: the cited results at $\mathbb G$; the structure conditions on $(A,H)$;
$\Om$ bounded,
open, with non-empty complement; $u,v$ continuous on $\bar\Om$; the two viscosity properties;
and $u\le v$ on $\del\Om$. Its two conclusions are
\begin{itemize}
\item if $v$ is constant on $\bar\Om$ then $u\le v$ in $\Om$;
\item if $u$ is constant on $\bar\Om$ then $u\le v$ in $\Om$.
\end{itemize}
Case (ii) supplies one of the two antecedents, so one of the two conclusions applies.

\medskip\noindent\textbf{Case (iii).} The hypotheses of \noderef{UncondCaseFirstLayer} are
again exactly the ambient hypotheses of the present statement, and this is the one case that
consumes both cited-results hypotheses: that node asks for a package at $\mathbb G$ and a
package at the Euclidean setting $\mathbb G_E$ on $\mathbb R^{m}$ (\noderef{EuclidSetting} at
$N:=m$), and the present statement carries exactly those two. Its two conclusions are
\begin{itemize}
\item if $r=1$ then $u\le v$ in $\Om$;
\item if there are $F,Gf$ continuous on $\pi_1(\bar\Om)$ with $u=F\circ\pi_1$ and
  $v=Gf\circ\pi_1$, then $u\le v$ in $\Om$.
\end{itemize}
Case (iii) supplies one of the two antecedents, so one of the two conclusions applies.

\medskip\noindent In each of the three cases we obtain $u(x)\le v(x)$ for every $x\in\Om$,
which is the assertion.
\end{proof}
